\documentclass[12pt,reqno]{amsart}
\usepackage{amssymb}
\usepackage{amsmath, mathtools}

\usepackage{amsthm}

\usepackage{amscd}

%     Roman numerals
\newcommand{\RNum}[1]{\uppercase\expandafter{\romannumeral #1\relax}}


\usepackage[T1]{fontenc} % Rendering compatibility for unavailable T2A; original Latin text unchanged.
\usepackage[utf8]{inputenc}
\usepackage[english]{babel}

\input epsf.tex

\usepackage{xcolor,colortbl}
\def\red {\color{red} }
\def\blue{\color{blue} }
\def\green{\color{green} }
\long\def\cyan{\color{cyan} }
\definecolor{darkblue}{rgb}{0.0, 0.0, 0.55}
\definecolor{darkcyan}{rgb}{0.0, 0.55, 0.55}
\definecolor{darkmagenta}{rgb}{0.55, 0.0, 0.55}
\definecolor{denim}{rgb}{0.08, 0.38, 0.74}
\definecolor{indigo}{rgb}{0.29, 0.0, 0.51} %% indigo(web)
\definecolor{lavenderindigo}{rgb}{0.58, 0.34, 0.92}
\definecolor{limegreen}{rgb}{0.2, 0.8, 0.2}
\definecolor{lincolngreen}{rgb}{0.11, 0.35, 0.02}
\definecolor{napiergreen}{rgb}{0.16, 0.5, 0.0}
\definecolor{LightCyan}{rgb}{0.88,1,1}
\definecolor{emerald}{rgb}{0.31, 0.78, 0.47}

\def\myitemmargin{
\leftmargini=28pt %originally 37.89pt
\leftmarginii=20pt %originally 23.93pt
\leftmarginiii=15pt %originally 30.13pt
\leftmarginiv=10pt  %originally 24.90
}

\newdimen\xsize
\newdimen\oldbaselineskip
\newdimen\oldlineskiplimit
\xsize=.7\hsize
\def\nolineskip{\oldbaselineskip=\baselineskip\baselineskip=0pt%
\oldlineskiplimit=\lineskiplimit\lineskiplimit=0pt}
\def\restorelineskip{\baselineskip=\oldbaselineskip%
\lineskiplimit=\oldlineskiplimit}
%
\def\putm[#1][#2]#3{% #1 X-axis  #2 Y-axis(downwords)  #3 $"mathobject"$
\hbox{\vbox to 0pt{\parindent=0pt%
\vskip#2\xsize\hbox to0pt{\hskip#1\xsize $#3$\hss}\vss}}}%
%\enddef\putm
%
%\def\putm[#1][#2]#3{% #1 X-axis  #2 Y-axis(downwords)  #3 $"mathobject"$
%\vbox to 0pt{\noindent\hskip#1\xsize\lower#2\xsize\hbox to0pt{$#3$\hss}\vss}}%
%\enddef\putm
%
\def\putt[#1][#2]#3{% #1 X-axis  #2 Y-axis(downwords)  #3 "text"
\vbox to 0pt{\noindent\hskip#1\xsize\lower#2\xsize%
\vtop{\restorelineskip#3}\vss}}
%\enddef\putt


%
\font\Bigbf=cmbx12 scaled 1200
\font\bigbf=cmbx12
%
\thinmuskip = 2mu
\medmuskip = 2.5mu plus 1.5mu minus 2.1mu  %minus 2.5mu
\thickmuskip = 4mu plus 6mu
%
%\font\teneusm=eusm10
%\font\seveneusm=eusm7
%\font\fiveeusm=eusm5
%\newfam\eusmfam
%\textfont\eusmfam=\teneusm
%\scriptfont\eusmfam=\seveneusm
%\scriptscriptfont\eusmfam=\fiveeusm
%\def\cal#1{{\fam\eusmfam\relax#1}}
%
%\def\cal#1{{\mathcal{#1}}}



\DeclareFontFamily{U}{rsf}{\skewchar\font'177}%
\DeclareFontShape{U}{rsf}{m}{n}{<-6>rsfs5<6-8>rsfs7<8->rsfs10}{}%
\DeclareFontShape{U}{rsf}{b}{n}{<-6>rsfs5<6-8>rsfs7<8->rsfs10}{}%
\DeclareMathAlphabet\RSFS{U}{rsf}{m}{n}
\SetMathAlphabet\RSFS{bold}{U}{rsf}{b}{n}
\DeclareOption{psamsfonts}{%
  \DeclareFontFamily{U}{rsf}{}%
  \DeclareFontShape{U}{rsf}{m}{n}{<-6>rsfs5<6-8>rsfs7<8->rsfs10}{}%
  \DeclareFontShape{U}{rsf}{b}{n}{<-6>rsfs5<6-8>rsfs7<8->rsfs10}{}%
}
\ProcessOptions

\def\rfs#1{{\RSFS{#1}}}
\let\scr=\rfs
%\let\scr=\mathcal 
%\def\scr#1{{\mathcal{#1}}}



%
\def\mib#1{\boldsymbol{#1}} \let\bsy=\mib
\def\sf#1{{\mathsf{#1}}} %\def\sf{\fam\sffam}
\def\bfsl#1{{\bfseries \slshape #1}}
\def\slsf#1{{\slshape \sffamily #1\/}}

%\let\rm= \rmfamily %roman
%\let\sf= \sffamily %sans serif
%\let\tt= \ttfamily %typewriter

%\let\rm =\mdseries %medium
%\let\bf =\bfseries %bold

%\let\rm= \upshape %upright
%\let\sl= \slshape %slanted
%\let\sl= \itshape %italic
%\let\sl= \scshape %small capital

%
\def\fr#1{{\mathfrak{#1}}}



% smalls
\def\msmall#1{\mathchoice{\hbox{\small$\displaystyle {#1}$}}{#1}{#1}{#1}}
\def\mSmall#1{\mathchoice{\hbox{\Small$\displaystyle {#1}$}}{#1}{#1}{#1}}
\def\smfrac#1#2{\msmall{\frac{#1}{#2}}}
\def\smatr#1{{\begin{smallmatrix} #1 \end{smallmatrix}}}
\def\spmatr#1{{\left(\begin{smallmatrix} #1 \end{smallmatrix}\right)}}

%
%
\hyphenation{Lip-schit-zian Lip-schitz com-pact-ness sub-mani-fold
 sub-mani-folds sub-fa-mily}
%
\let\3=\ss


\def\sfa{{\sf{A}}}
\def\ab{{\sf{ab}}}
\def\arg{{\sf{Arg}}}
\def\arf{{\sf{Arf}}}
\def\Alt{{\sf{Alt}}}
\def\aut{{\sf{Aut}}}
   \def\inn{{\sf{Inn}}}
   \def\Ad{{\sf{Ad}}}
   \def\out{{\sf{Out}}}
   \def\conj{{\sf{Conj}}}
\def\loc{{\sf{loc}}}
\def\st{{\mathsf{st}}}
\def\topo{{\mathsf{top}}}
%\def\st{\mathchoice{_\sf{st}}{_\sf{st}}%
%{_{\hbox{\fivesf st}}}{_{\hbox{\fivesf st}}}}
\def\area{\sf{area}}
\def\sfb{\sf{b}}
\def\br{\sf{Br}}
\def\sfpb{{\sf{PBr}}}
\def\pbr{\sf{PBr}}
%\def\fr{\sf{Fr}}
\def\const{{\sf{const}}}
\def\acosh{\sf{acosh}\,}
\def\coker{\sf{Coker}\,}
\def\codim{\sf{codim}\,}
\def\divi{{\sf{Div}\,}}
\def\dscr{\sf{Dscr}}
\def\deg{\sf{deg}\,}
\def\diam{\sf{diam}\,}
\def\dim{\sf{dim}\,}
\def\dimc{\dim_\cc}
\def\dimr{\dim_\rr}
\def\Def{\sf{Def}}
\def\dfrm{\sf{dfrm}}
\def\dist{\sf{dist}\,}
\def\endo{\sf{End}}
\def\endr{\sf{End}\vph_\rr}
\def\endrn{\endr(\cc^n)}
\def\endc{\sf{End}\vph_\cc}
\def\es{{\sf{es}}}
\def\ev{\sf{ev}}
\def\ext{{\sf{Ext}}}
\def\fin{\sf{Fin}}
%\def\bff{\mib{F}}
\def\gh{\sf{gh}}
\def\ind{\sf{ind}}
\def\ini{\sf{Ini}}
\def\integ{\sf{int}}
\def\length{\sf{length}}
\def\sym{{\sf{Sym}}}
\def\sind{\sf{S\hbox{-}ind}\,}
\def\gcd{\sf{gcd}}
\def\gr{{\sf{Gr}}}
\def\frg{{\fr{g}}}
\def\frh{{\fr{h}}}
\def\sfh{{\sf{H}}}
\def\smsfh{\sf{h}}
\def\hom{\sf{Hom}\,}
\def\scrhom{{\scrh\!\!o\!m}}
\def\homr{\sf{Hom}\vph_\rr}
\def\hur{{\mib{H}}}%{\fr{H}}
\def\id{\sf{Id}}
\def\im{\sf{Im}\,}
\def\bint{{-}\mathchoice{\mkern-18mu}{\mkern-15mu}{\mkern-15mu}{\mkern-15mu}\int}
\def\sfk{{\sf{K}}}
\def\sfl{{\sf{L}}}
\def\re{\sf{Re}\,}
\def\ker{\sf{Ker}\,}
\def\length{\sf{length}}
\def\mod{\sf{mod}}
\def\sfm{{\sf{M}}}  %% \def\sfm{\sf{m}}
\def\bfm{{\boldsymbol{m}}}
\def\nod{{\sf{nod}}}
\def\ord{\sf{ord}}
\def\sord{\sf{s\text{-}ord}}
\def\osc{\sf{osc}}
\def\mosc{\textsf{M{-}osc}}
\def\sfp{{\sf{P}}}
\def\pr{\sf{pr}}
\def\pt{\sf{pt}}
\def\rank{\sf{rank}}
%\def\reg{\sf{reg}}   
\def\res{\sf{Res}}
\def\pic{\sf{Pic}}
\def\Sing{{\sf{Sing}}}%\def\sing{\sf{sing}}
\def\Sl{{\mib{Sl}}}
\def\frsl{{\frak{sl}}}
\def\frso{{\frak{so}}}
\def\fro{{\frak{o}}}
\def\SO{{\mib{SO}}}
\def\SU{{\mib{SU}}}
\def\UU{{\mib{U}}}
\def\Sp{{\mib{Sp}}}
\def\sp{{\fr{sp}}}
\def\pgl{\mib{PGl}}
\def\Gl{\mib{Gl}}
\def\Aff{\mib{A\!f\!\!f}}
\def\uone{{\mib{U}(1)}}
%\def\spec{\sf{Sp}}
\def\frs{{\fr s}}
\def\fraks{{\fr S}}
\def\tan{\sf{tan}}
\def\tr{{\sf{tr}}}
\def\tor{{\sf{tor}}}
\def\vol{\sf{Vol}}
\def\sfw{\sf{W}}
\def\sfz{\sf{Z}}


\def\bfone{\boldsymbol{1}}
\def\mbfc{{\mib{c}}}
\def\mbfe{{\mib{e}}}
\def\mbff{{\mib{f}}}
\def\mbfj{{\mib{j}}}
\def\mbfk{{\mib{k}}}
\def\mbfl{{\mib{l}}}
\def\mbfm{{\mib{m}}}
\def\mbfn{{\mib{n}}}
\def\mbfp{{\mib{p}}}
\def\mbfs{{\mib{s}}}
\def\mbft{{\mib{t}}}
\def\mbfu{{\mib{u}}}
\def\mbfv{{\mib{v}}}
\def\mbfw{{\mib{w}}}
\def\mbfx{{\mib{x}}}
\def\mbfy{{\mib{y}}}
\def\mbfz{{\mib{z}}}


\def\bfa{{\mib{A}}}
\def\bfb{{\mib{B}}}
\def\bfc{{\mib{C}}}
\def\bfd{{\mib{D}}}
\def\bfds{{\mib{DS}}}
\def\bfe{{\mib{E}}}
\def\bff{{\mib{F}}}
\def\bfh{{\mib{H}}}
\def\bfl{{\mib{L}}}
\def\bfn{{\mib{N}}}
\def\bfo{{\mib{O}}}
\def\bfp{{\mib{P}}}
\def\bfq{{\mib{Q}}}
\def\bfr{{\mib{R}}}
\def\bfs{{\mib{S}}}
\def\bft{{\mib{T}}}
\def\bfu{{\mib{U}}}
\def\bfv{{\mib{V}}}
\def\bfy{{\mib{Y}}}

\def\vpp{\,\,\vec{\vphantom{p}}\!\!p'{}}
\def\vdp{\,\,\vec{\vphantom{d}}\!\!d'{}}

\def\dash{{\text{--}}}

\let\del=\partial
\let\bbx=\blacksquare
\def\eps{\varepsilon}
\def\epsi{\varepsilon}
\def\dprime{{\prime\prime}}
\def\prist{{\prime*}}
\def\dprist{{\prime\prime*}}
\let\vkappa=\varkappa
\let\brel=\buildrel
\def\<{\langle}\let\lan=\<
\def\>{\rangle}\let\ran=\>
\def\bss{\backslash} \let\bs=\bss
\def\bsq{{\mathchoice{\rule{3.5pt}{3.5pt}}%
{\rule{3.5pt}{3.5pt}}%
{\rule{3.pt}{3.pt}}%
{\rule{2.2pt}{2.2pt}}%
}}
\def\comp{\Subset}
\def\d{\partial}
\def\dbar{{\barr\partial}}
\def\barz{{\bar z}}
\def\dbarj{\dbar_J}
\def\ddef{\mathrel{{=}\raise0.3pt\hbox{:}}}
\def\deff{\mathrel{\raise0.3pt\hbox{\rm:}{=}}}
\def\inv{^{-1}}
\def\hook{\hookrightarrow}
\def\fraction#1/#2{\mathchoice{{\msmall{ \frac{#1}{#2} }}}%
{{ \frac{#1}{#2} }}{{#1/#2}}{{#1/#2}}}
%\def\norm#1{\Vert #1 \Vert} % Roman defined his own \norm
\def\bnorm#1{\left\Vert #1 \right\Vert}
\def\half{{\fraction1/2}}
\def\le{\leqslant}
\def\leq{\leqslant}
\def\vph{^{\mathstrut}}
\def\vpht{^{\vphantom{2}}}
\def\lrar{\longrightarrow}
\def\llar{\longleftarrow}
\def\xarr#1{\xrightarrow{#1}}
\def\emptyset{\varnothing}
\def\scirc{\mathop{\mathchoice{\hbox{\small$\circ$}}{\hbox{\small$\circ$}}%
{{\scriptscriptstyle\circ}}{{\scriptscriptstyle\circ}}}}
%
\def\longpoints{\leaders\hbox to 0.5em{\hss.\hss}\hfill \hskip0pt}
\def\stateskip{\smallskip}
\def\state#1. {\stateskip\noindent{\it#1.}} %\medskip
\def\statebf#1. {\stateskip\noindent{\bf#1.}} %\medskip
\def\statep#1. {\stateskip\noindent{\bf#1}} %\medskip

\def\step#1{{\slsf{Step \ref{#1}}}}  %\def\step#1{{\sl Step #1)}}
\def\Chi{\raise 2pt\hbox{$\chi$}}
\def\eg{\,{\sl{e.g.}}}
\def\ie{\,{\sl i.e.}}
%\def\iff{if and only if }
%{\hskip1pt plus1pt{\sl iff\/\hskip2pt plus1pt }}
\def\wrt{with respect to }
\def\isl{{\text{\slsf{i}}\mskip1mu}} %{{{\text{\bf i$\mskip1mu$}}}}
\let\cpi=\isl
\def\sli{{\sl i)\/} \,}         \def\slip{{\sl i)}}
\def\slii{{\sl i$\!$i)\/} \,}   \def\sliip{{\sl i$\!$i)}}
\def\sliii{{\sl i$\!$i$\!$i)\/} \,}    \def\sliiip{{\sl i$\!$i$\!$i)}}
\def\sliv{{\sl i$\!$v)\/} \,}   \def\slivp{{\sl i$\!$v)}}
\def\slv{{\sl v)\/} \,}         \def\slvp{{\sl v)}}
\def\slvi{{\sl v$\!$i)\/} \,}   \def\slvip{{\sl v$\!$i)}}
\def\slvii{{\sl v$\!$i$\!$i)\/} \,}   \def\slviip{{\sl v$\!$i$\!$i)}}
\def\jac{{\sf{Jac}}}
\def\reg{{\sf{reg}}}
\def\sing{{\sf{sing}}}
\def\mat{\sf{Mat}}

\def\sla{{\sl a)\/} \,}
\def\slb{{\sl b)\/} \,}
\def\slc{{\sl c)\/} \,}

\def\fol{\sf{Fol}}
\def\gr{\sf{Gr}}
\def\pd{\sf{PD}}

\def\matn{\mat(n,\cc)}
\def\det{\sf{det}}
\def\diff{\text{\normalfont{Diff}}}
\def\aut{\text{\normalfont{Aut}}}
\def\difff{{\diff\mskip 4mu{}^!\mskip-2mu}}
\def\map{{\mathsf{Map}}}
\def\smap{{\textsf{S-Map}}}
\def\whmap{\wh\map\vphantom{M}}
\def\wtmap{\wt\map\vphantom{M}}
\def\maph{\map^\rmh}
%\def\map{{{\scr M}\mskip -5mu a\mskip -1.5mu p}}
\def\star{\mathop{\msmall{*}}}
\def\barr#1{\mskip1mu\overline{\mskip-1mu{#1}\mskip-1mu}\mskip1mu}
%\def\barr#1{\overline{#1}\vphantom{#1}}
%\def\usig{U\!\bfsig}
\def\symp{\text{\normalfont{Symp}}}
\let\sympl=\symp
\def\ham{\text{\normalfont{Ham}}}
\def\cmp{{\sf{comp}}}
\def\can{{\sf{can}}}
\def\tic{{\ti c}}
\def\tiw{{\ti w}}
\def\tix{{\,\wt{\times}\,}}
%
\def\Chi{\raise 2pt\hbox{$\chi$}}
\def\yps{\Upsilon}
%\let\vphi=\phi\let\phi=\varphi\let\varphi=\vphi
%
\def\igam{{\mathchar"0100}} % Italic \Gamma
\def\idel{{\mathchar"0101}} % Italic \Delta
\def\ithe{{\mathchar"0102}} % Italic \Theta
\def\ilam{{\mathchar"0103}} % Italic \Lambda
\def\ixi{{\mathchar"0104}} % Italic \Xi
\def\ipi{{\mathchar"0105}} % Italic \Pi
\def\isig{{\mathchar"0106}} % Italic \Sigma
\def\iyp{{\mathchar"0107}} % Italic \Upsilon
\def\iphi{{\mathchar"0108}} % Italic \Phi
\def\ipsi{{\mathchar"0109}} % Italic \Psi
\def\iom{{\mathchar"010A}} % Italic \Omega


%
\def\bfalpha{{\boldsymbol\alpha}}%
\def\bfbeta{{\boldsymbol\beta}}%
\def\bfgamma{{\boldsymbol\gamma}}\let\bfga=\bfgamma%
\def\bfdelta{{\boldsymbol\delta}}\let\bfde=\bfdelta%
\def\bfepslon{{\boldsymbol\varepsilon}}%
\def\bfvarepsilon{{\boldsymbol\epslon}}%
\def\bfzeta{{\boldsymbol\zeta}}%
\def\bfeta{{\boldsymbol\eta}}%
\def\bftheta{{\boldsymbol\theta}}\let\bfth=\bftheta%
\def\bfvartheta{{\boldsymbol\vartheta}}%
\def\bfiota{{\boldsymbol\iota}}%
\def\bfkappa{{\boldsymbol\varkappa}}%
\def\bfvarkappa{{\boldsymbol\kappa}}%
\def\bflambda{{\boldsymbol\lambda}} 
   \let\bflamb=\bflambda%
   \let\bflam=\bflambda%
   \let\bfla=\bflambda%
\def\bfmu{{\boldsymbol\mu}}%
\def\bfnu{{\boldsymbol\nu}}%
\def\bfxi{{\boldsymbol\xi}}%
\def\bfpi{{\boldsymbol\pi}}%
\def\bfvarpi{{\boldsymbol\varpi}}%
\def\bfrho{{\boldsymbol\rho}}%
\def\bfvarrho{{\boldsymbol\varrho}}%
\def\bfsigma{{\boldsymbol{\sigma}}}
    \def\bfsig{\bfsigma}%
\def\bfSigma{{\boldsymbol{\Sigma}}}
    \def\bfSig{\bfSigma}%
\def\bfvarsigma{{\boldsymbol\varsigma}}%
\def\bftau{{\boldsymbol\tau}}%
\def\bfupsilon{{\boldsymbol\upsilon}}%
\def\bfUps{{\boldsymbol\Upsilon}}%
\def\bfphi{{\boldsymbol\varphi}}%
\def\bfvarphi{{\boldsymbol\phi}}%
   \def\bffi{{\bfvarphi}}
\def\bfchi{{\boldsymbol\chi}}%
\def\bfpsi{{\boldsymbol\psi}}%
\def\bfPsi{{\boldsymbol\Psi}}%
\def\bfomega{{\boldsymbol\omega}}%
    \let\bfom=\bfomega
\def\bfell{{\boldsymbol\ell}}%
\def\bfz{{\boldsymbol z}}%
%
\def\scra{\scr{A}}
  \def\cala{{\mathcal A}}
  \def\rma{{\mathrm A}}
  \def\frA{{\mathfrak{A}}}
\def\scrb{\scr{B}}
  \def\calb{{\mathcal B}}
   \def\rmb{{\mathrm{B}}}
   \def\rmbd{{\mathrm{BD}}}
   \def\rmbe{{\mathrm{BE}}}
\def\scrc{\scr{C}}
  \def\calc{{\mathcal C}}
\def\scrd{\scr{D}}
   \def\rmd{{\mathrm{D}}}
  \def\cald{{\mathcal D}}
\def\scre{\scr{E}}
   \def\cale{{\mathcal E}}
   \def\rme{{\mathrm{E}}}
\def\scrf{\scr{F}}
  \def\calf{{\mathcal F}}
\def\scrg{\scr{G}}
  \def\calg{{\mathcal G}}
\def\scrh{\scr{H}}
  \def\calh{{\mathcal H}}
  \def\rmh{{\mathrm{H}}}
  \def\diffh{\diff^{\;\rmh}}
\def\scri{\scr{I}}
  \def\cali{{\mathcal I}}
  \def\IW{{\sfw\cali}}
  \def\rmi{{\mathrm{I}}}
\def\scrj{\scr{J}}
  \def\calj{\mathcal{J}}
\def\scrk{\scr{K}}
  \def\calk{{\mathcal K}}
\def\scrl{\scr{L}}
  \def\call{{\mathcal L}}
  \def\rml{{\mathrm{L}}}
\def\scrm{\scr{M}}
  \def\whcalm{\wh{\scr M}}
  \def\barm{{\mskip5mu\barr{\mskip-5mu\scr M}}{}}
    \def\calm{{\mathcal M}}
% \def\frm{{\fr{M}}} %!%! Used by XY-pic
\def\scrn{{\scr N}}
    \def\caln{{\mathcal N}}
  \def\rmn{{\mathrm{N}}}
\def\scro{{\scr O}}
\def\scrp{{\scr P}}
\def\scrpp{{\scr{P\!\!P}}}
  \def\calp{{\mathcal P}}
%  \def\whcalp{\wh{\scr P}}
\def\scrq{{\scr Q}}
  \def\calo{{\mathcal O}}
\def\scrr{{\scr R}}
  \def\calr{{\mathcal R}}
\def\scrs{{\scr S}}
  \def\cals{{\mathcal S}}
  \def\osa{{\iom\cals\cala}}
  \def\osj{{\iom\cals\calj}}
  \def\ohj{{\iom\calj}}
\def\scrt{{\scr T}}%{{\mathcal T}}
  \def\calt{{\mathcal T}}
\def\scru{{\scr U}}
  \def\calu{{\mathcal U}}
\def\scrv{{\scr V}}
  \def\calv{{\mathcal V}}
\def\scrw{{\scr W}}
  \def\calw{{\mathcal W}}
\def\scrx{{\scr X}}
  \def\calx{{\mathcal X}}
\def\scry{{\scr Y}}
  \def\caly{{\mathcal Y}}
\def\scrz{{\scr Z}}
  \def\calz{{\mathcal Z}}


\let\warr=\rightsquigarrow
\def\hrar{\hookrightarrow}
\let\xrar=\xrightarrow
\let\xlar=\xleftarrow
\def\xlrar#1{{\leftarrow\mskip-11mu\xrar{#1}}} 
\def\xmapsto#1{\mapstochar\xrar{#1}}
\def\trans{\pitchfork}
\def\eps{\varepsilon}
\def\epsi{\varepsilon}
\def\bs{\backslash}
\def\ogran{{\hskip0.7pt\vrule height6pt depth2.5pt\hskip0.7pt}}
\def\comp{\Subset}
\def\d{\partial}
\def\dbar{{\barr\partial}}
\def\1{{1\mkern-5mu{\rom l}}}
\def\dbarj{\dbar_J}
\def\ge{\geqslant}\def\geq{\geqslant}
\def\inv{^{-1}}
\let\wh=\widehat
\let\wt=\widetilde
%\def\fraction#1/#2{\mathchoice{{\msmall{ #1\over#2}}}%
%{{ #1\over #2 }}{{#1/#2}}{{#1/#2}}}
\def\half{{\fraction1/2}}
\def\le{\leqslant}
\def\vph{^{\mathstrut}}
\def\emptyset{\varnothing}
\def\parnin{\par\noindent}

\def\ti#1{{\tilde{#1}}}
\def\term#1#2{\underline{#2}_{[{#1}]}}

\def\comment#1\endcomment{}

\abovedisplayskip=5pt plus3pt minus2pt
 \belowdisplayskip=\abovedisplayskip

%\show\abovedisplaylimit
%\def\o#1{\buildrel\circ \over{#1}}$\o u$


%
%%%

\def\lineeqqno(#1){\hfill\llap{\vbox to 10pt%
{\vss\begin{align} \eqqno(#1)\end{align}\vss}}\vskip1pt}

% \def\Matrix{\begin{matrix}}
% \def\Endmatrix{\end{matrix}}

\def\pmatr#1{{\left(\begin{matrix} #1 \end{matrix} \right)}}
\def\spmatr#1{{\left(\begin{smallmatrix} #1 \end{smallmatrix}\right)}}

% \def\Cases{\begin{cases}}
% \def\Endcases{\end{cases}}


\textwidth= 6.3 truein
\textheight=9.2 truein
\vsize=9.3 truein
\voffset= -.5truein
\oddsidemargin 0.1in
\evensidemargin 0in
\advance\headheight 1.2pt

\newtheorem{theorem}{Theorem}[section]
\newtheorem{proposition}[theorem]{Proposition}
\newtheorem{conjecture}[theorem]{Conjecture}
\newtheorem{problem}[theorem]{Problem}

\newtheorem*{theorem*}{Theorem}
\newtheorem{lemma}[theorem]{Lemma}
\newtheorem{exercise}[theorem]{Exercise}
\newtheorem{corollary}[theorem]{Corollary}
\newtheorem{definition}[theorem]{Definition}

\newtheorem{example}[theorem]{Example}

\theoremstyle{remark}
\newtheorem{remark}[theorem]{Remark}

\newcounter{stepnr}
\newcounter{laststepnr}
\newenvironment{stepping}%
{\begin{list}%
{ \slsf{Step \arabic{stepnr}.} }%
{\usecounter{stepnr}  \setlength{\leftmargin}{0pt} 
\setlength{\itemindent}{35pt} \setlength{\listparindent}{0pt} 
\setlength{\itemsep}{7pt plus 3pt}
}}%
{\end{list}}

\def\reffig#1{{\slsf{Figure \ref{#1}}}}



%\numberwithin{equation}{section}
%\renewcommand{\theequation}{\thesection.\arabic{equation}.{[\thelabel]}}
%\renewcommand{\theequation}{\thesection\thesubsection\arabic{equation}}

%old \eqq...
%\def\eqqno(#1){\label{eq#1}%
%      \index{\slsf{ Eqtn \ref{eq#1}},  \ label {\tt(#1)}|Zpage}}
%\def\eqqref(#1){\eqref{eq#1}}
%\def\eeq{\end{equation}}

%%% Local Variables: 
%%% mode: latex
%%% TeX-master: "invol.tex"
%%% End: 
 %Gleb's definitions

\usepackage[sort]{cite}
\usepackage{tikz-cd}
\usetikzlibrary{cd}
%\usepackage{mdwlist, enumitem}
\usepackage{dirtytalk}
\usepackage[linktoc=page, colorlinks, linkcolor=blue, citecolor=blue]{hyperref}

\usepackage{xcolor}
\usepackage{centernot}

\usepackage{enumitem}

\usepackage{pgfplots}
\usepackage{multicol}

\pgfplotsset{compat=1.17}

\makeatletter
\renewenvironment{proof}[1][\proofname]{%
  \par\vspace{\topsep}%
  \normalfont\topsep6\p@\@plus6\p@\relax
  \trivlist
  \item[\hskip\labelsep\itshape #1\@addpunct{.}]\ignorespaces
}{%
  \endtrivlist
}
\makeatother

\numberwithin{equation}{section}



%------Roman's definitions---------
\DeclarePairedDelimiter \abs{\lvert}{\rvert} % \abs{...}, \abs*{...}, \abs[\Big]{...}, etc.
\DeclarePairedDelimiter \norm{\lVert}{\rVert}
\DeclarePairedDelimiterX \ip[2]{\langle}{\rangle}{#1,#2}
\DeclarePairedDelimiterXPP \Prob[1]{\mathbb{P}}\{\}{}{\newcommand\given{\nonscript\:\delimsize\vert\nonscript\:\mathopen{}} #1} % \Prob{...}, \Prob*{..}, \Prob[\Big]{...}, \Prob{A \given B}
\DeclarePairedDelimiterXPP \Probevent[1]{\mathbb{P}}(){}{\newcommand\given{\nonscript\:\delimsize\vert\nonscript\:\mathopen{}}#1} % \Probevent{...}, \Prob*{..}, \Probevent[\Big]{...}, \Probevent{A \given B}

\DeclareMathOperator{\disc}{disc}
\DeclareMathOperator{\E}{\mathbb{E}}

\newcommand{\one}{\mathbf{1}}

\def \R {\mathbb{R}}
\def \e {\varepsilon}
\def \l {\lambda}

%----------------------------------


\parindent = 0pt

\usepackage[mathcal]{euscript}

\usepackage{titlesec}
\titleformat{\section}[runin]{\bfseries}{\thesection.}{3pt}{}[.]


%\myitemmargin 
\baselineskip =15.0pt plus 2.5pt
\usepackage{geometry}
\newgeometry{vmargin={25mm}, hmargin={22mm,22mm}, footskip=10mm}   % set the margins




\begin{document}

\title[Optimal discrepancy for triangles]{Optimal discrepancy for triangles}

\author{Gleb Smirnov}
\address{
Mathematical Sciences Institute, 
Australian National University, Canberra, Australia}
\email{gleb.smirnov@anu.edu.au}


%\dedicatory{}

\begin{abstract}
We prove that every $n$-point set in the plane has a two-coloring with discrepancy $O(n^{1/4})$ for all triangles. This improves the previous bound $O(n^{1/4}\sqrt{\log n})$ and matches the classical lower bound.
\end{abstract}

\maketitle

\setcounter{section}{0}

%\section{Main results} \label{intro}

\section{Main results}\label{intro}
Let \(X=\{p_1,\ldots,p_n\}\subset\R^2\) be a set of points. We color each point either \(+1\) or \(-1\). We want every triangle to contain nearly as many \(+1\)'s as \(-1\)'s. To make this precise, encode a coloring by a vector \(\chi\in\{-1,1\}^n\): the point \(p_i\) gets color \(\chi_i\). For a subset 
\(F\subset X\), write:
\[
\chi(F):=\sum_{p_i\in F}\chi_i.
\]
We want \(|\chi(X\cap T)|\) to be small for all triangles \(T\) at once.
\smallskip%

More generally, let \(\mathcal R\) be a family of subsets of \(\R^2\). The \emph{discrepancy} of \(X\) wrt \(\mathcal R\) is the imbalance of the best coloring:
\[
\disc(X,\mathcal R):=\min_{\chi\in\{-1,1\}^n}\ \sup_{R\in\mathcal R}\,|\chi(X\cap R)|.
\]
Our main result is the following statement.
\begin{theorem}\label{thm:triangles}
Every set \(X=\{p_1,\ldots,p_n\}\subset\R^2\) has a coloring \(\chi\in\{-1,1\}^n\) such that:
\[
|\chi(X\cap T)|\le C\,n^{1/4}\qquad\text{for all triangles }T,
\]
where \(C>0\) is an absolute constant.
\end{theorem}
For planar triangles, the best bound we are aware of is \(O(n^{1/4}\sqrt{\log n})\). This bound follows from more general results of Matou\v{s}ek et al.\ \cite{MWW}. See Chazelle \cite[\S\,1.4]{Chaz} for a simple exposition. 
\smallskip%

Theorem~\ref{thm:triangles} is sharp. The lower bound is a classical result of Beck \cite{Beck-2}. A rectangle is a union of two triangles. Hence, it suffices to prove the lower bound for rectangles in arbitrary position. Chazelle \cite[\S\,3.2]{Chaz} and Beck \cite{Beck-2} explain the following lower bound for their Lebesgue discrepancy. For every \(n\)-point set \(X\subset[0,1]^2\), there is a rotated rectangle \(R\subset[0,1]^2\) such that:
\[
\bigl|\#(X\cap R)-n\operatorname{area}(R)\bigr|=\Omega(n^{1/4}).
\]
Beck's transference principle \cite{Beck-1} carries this bound over to the combinatorial discrepancy.
\smallskip%

We deduce Theorem~\ref{thm:triangles} from a statement about halfplanes. A \emph{halfplane trace} is a set \(X\cap P\), where \(P\subset\R^2\) is a halfplane. Let \(\mathcal H(X)\) denote the family of halfplane traces. Let the points of \(X\) be ordered. For \(0\le k\le n\), set:
\[
X_k:=\{p_1,\ldots,p_k\}.
\]
\begin{theorem}\label{thm:prefix1}
Every ordered \(X=\{p_1,\ldots,p_n\}\subset\R^2\) has a coloring \(\chi\in\{-1,1\}^n\) such that:
\[
|\chi(F\cap X_k)|\le C\,n^{1/4}\qquad\forall F\in\mathcal H(X),\ \ 0\le k\le n,
\]
where \(C>0\) is an absolute constant.
\end{theorem}
For \(k=n\), this is Matou\v{s}ek's bound for halfplanes \cite{Mat-halfspaces}. The point is that one coloring controls all prefixes at once. We explain the step from halfplanes to triangles in \S\,\ref{sec:triangles}.
\smallskip%

We prove Theorem~\ref{thm:prefix1} using the Fisher information technique of \cite{Versh-Sm-3, Versh-Sm-4}. The halfplane constraints cut out a convex body in \(\R^n\); it is an intersection of slabs. The crucial step is to bound the Fisher information of densities supported on intersections of slabs. We derive such bounds in \S\,\ref{sec:slabs}; they may be of independent interest. The full proof occupies \S\,\ref{sec:slabs}, \S\,\ref{chaining}, and \S\,\ref{sec:prefix}.
\smallskip%

We now discuss other related work.
\smallskip%

Theorem~\ref{thm:prefix1} is a vector balancing statement. Indeed, for each \(p_i\), let \(u_i\in\{0,1\}^{\mathcal H(X)}\) be the indicator of the halfplane traces that contain \(p_i\). Then Theorem~\ref{thm:prefix1} balances all prefix sums \(\sum_{i\le k}\chi_iu_i\) in \(\ell_\infty\).
\smallskip%

Vector balancing has seen rapid progress. For the Koml\'os problem, Bansal and Jiang \cite{BansalJiang} improved Banaszczyk's \(O(\log^{1/2}n)\) bound to \(O(\log^{1/4}n)\). Guo, Fang and Lu \cite{GFL} obtained a constant bound; this settles the Koml\'os conjecture. See the work of Karingula and Lovett \cite{KarLov} and of Akbas and Sra \cite{AkbasSra-1, AkbasSra-2} for alternative proofs. In the online Koml\'os problem, the vectors arrive one by one, and all prefix sums must stay small. Optimal bounds are obtained by Kulkarni, Reis and Rothvoss \cite{Kulk-Reis-Roth} and by Aden-Ali \cite{AdenAli}. Their work improves earlier bounds of Alweiss, Liu and Sawhney \cite{AL-Saw} and Bansal et al.\ \cite{Bans-JMSS}. The online Spencer problem, where the vectors are only bounded in \(\ell_\infty\), was solved by Altschuler and Tikhomirov \cite{AltTikh-1}.
\smallskip%

Our proof relies on the partial coloring method. The method goes back to Beck \cite{Beck-3} and Spencer \cite{Spencer-6}. Constructive methods for partial coloring are due to Bansal \cite{Bansal-10} and Lovett and Meka \cite{LovMeka}. Our mechanism for partial coloring is different from all these. These methods need the convex body to be large: in entropy or in Gaussian measure. We need instead a density supported in it with small Fisher information.
\smallskip%

For discrepancy of geometric set systems, combinatorial and Lebesgue, see the monographs \cite{Mat, Chaz, Beck-Chen} and references therein. For halfplanes, the bound \(O(n^{1/4})\) is due to Matou\v{s}ek \cite{Mat-halfspaces}. Similar to that work, our proof relies on chaining and uses the packing lemma for halfplanes; see Mustafa~\cite{Mustafa} for a simple proof. For axis-parallel boxes, see Matou\v{s}ek et al. \cite{MNT} and 
Nikolov \cite{Nikolov}. On the Lebesgue discrepancy, see also Chen and Travaglini~\cite{ChenTrav} for convex polygons and Bilyk and Mastrianni~\cite{BilykMast} for rectangles with restricted rotations. By Beck's transference principle \cite{Beck-1}, an upper bound on combinatorial discrepancy gives a corresponding upper bound for Lebesgue discrepancy. In particular, Theorem~\ref{thm:triangles} gives \(O(n^{1/4})\) for the Lebesgue discrepancy of rotated rectangles. This improves Beck's bound \(O(n^{1/4}\sqrt{\log n})\) \cite{Beck-Chen}, and matches his lower bound.
\smallskip%

\textbf{Acknowledgment and AI assistance.} 
The author thanks OpenAI's ChatGPT for useful insights: ChatGPT suggested a simpler proof of Proposition~\ref{slabs}, while the original proof relied on the OU process. ChatGPT also helped the author with estimates in the proof of Proposition~\ref{prop:halfplanes}. Finally, ChatGPT helped to edit the manuscript.

\section{Fisher cost of intersecting slabs}\label{sec:slabs}

Let \(K\subset\R^n\) be a convex body. We seek a function \(\psi\in C^\infty(\R^n)\) with
\[
\|\psi\|_{L^2(\R^n)}=1,\qquad
\operatorname{supp}\psi\subset K,
\]
and with small
\[
\int_{\R^n}\|\nabla\psi\|_2^2.
\]
This quantity controls the Metropolis rejections below. In \cite{Versh-Sm-3}, it equals \(\operatorname{tr}M_\psi\). The bodies \(K\) we need are intersections of slabs. Here, we estimate \(\int\|\nabla\psi\|_2^2\) for such intersections.
\smallskip%

To begin with, we introduce a cutoff function. 
Fix an even function \(g\in C^\infty(\R)\), non-increasing on \([0,\infty)\), such that:
\[
g(t)=\begin{cases}
1, & |t|\le 1/2,\\
>0, & |t|<1,\\
0, & |t|\ge 1,
\end{cases}
\]
and such that \(\log g\) is \emph{concave} 
on \((-1,1)\). In particular, \(0\le g\le 1\). Set:
\[
U:=-2\log g\quad\text{on }(-1,1).
\]
Then \(U\) is convex, and \(\ddot U=0\) on \((-1/2,1/2)\). From
\[
g^2\ddot U=2\dot g^2-2g\ddot g,
\]
we get:
\begin{equation}\label{eq:gU}
\sup_{|t|<1}g(t)^2\,\ddot U(t)\le C,
\end{equation}
for some \(C\) depending only on \(g\).

\begin{proposition}\label{slabs}
Let \(v_1,\ldots,v_M\in\R^n\) and \(a_1,\ldots,a_M>0\). Set:
\[
K:=\bigl\{z\in\R^n:\ |\langle v_j,z\rangle|\le a_j\quad \forall j\bigr\},
\]
and assume \(K\) is bounded. Define:
\[
\psi(z):=Z^{-1/2}\,e^{-\|z\|_2^2/4}\prod_{j=1}^M
g\Bigl(\frac{\langle v_j,z\rangle}{a_j}\Bigr),
\]
where \(Z\) is chosen so that \(\|\psi\|_{L^2(\R^n)}=1\).
Then:
\begin{equation}\label{eq:slab-fisher}
4\int_{\R^n}\|\nabla\psi\|_2^2
\le
n+2\,C\sum_{j=1}^M\frac{\|v_j\|_2^2}{a_j^2}
\exp\Bigl(-\frac{a_j^2}{8\|v_j\|_2^2}\Bigr).
\end{equation}
\end{proposition}

\begin{proof}
Since \(0 \le g \le 1\), 
\[
Z\psi(z)^2\le e^{-\|z\|_2^2/2},
\]
with equality near \(z=0\). Hence, \(0<Z<\infty\). Clearly, \(\psi\in C^\infty(\R^n)\) and \(\operatorname{supp}\psi\subset K\). Since \(K\) is bounded, \(\psi\in C_c^\infty(\R^n)\).
\smallskip%

Let \(\Omega:=\{\psi>0\}\); this is the interior of \(K\). Write:
\[
\psi^2=Z^{-1}e^{-V},\quad
\text{where }\ V(z)=\frac{\|z\|_2^2}{2}+\sum_{j=1}^M U\Bigl(\frac{\langle v_j,z\rangle}{a_j}\Bigr).
\]
Then:
\begin{equation}\label{DeltaV}
\Delta V=n+\sum_{j=1}^M\frac{\|v_j\|_2^2}{a_j^2}\,\ddot U\Bigl(\frac{\langle v_j,z\rangle}{a_j}\Bigr)\quad\text{on }\Omega.
\end{equation}
Also, \(V=-2\log\psi-\log Z\) on \(\Omega\). Hence:
\[
\psi^2\Delta V=2\|\nabla\psi\|_2^2-2\psi\Delta\psi\quad\text{on }\Omega.
\]
Consequently,
\[
\int_\Omega\psi^2\Delta V=
\int_\Omega\bigl(2\|\nabla\psi\|_2^2-2\psi\Delta\psi\bigr).
\]
\(\psi\) vanishes outside \(\Omega\), and so does 
\(\nabla\psi\). Therefore:
\[
\int_\Omega\psi^2\Delta V=
\int_{\R^n}\bigl(2\|\nabla\psi\|_2^2-2\psi\Delta\psi\bigr).
\]
Since \(\psi\in C_c^\infty(\R^n)\), integration by parts gives:
\[
-\int_{\R^n}\psi\Delta\psi=\int_{\R^n}\|\nabla\psi\|_2^2.
\]
Consequently,
\[
\int_\Omega\psi^2\Delta V=4\int_{\R^n}\|\nabla\psi\|_2^2.
\]
Let us substitute \eqref{DeltaV} for \(\Delta V\). 
By \eqref{eq:gU}, each \(\ddot U(\langle v_j,z\rangle/a_j)\,\psi^2\) is at most \(C/Z\); and \(\Omega\) is bounded. So each term is integrable, and we obtain:
\begin{equation}\label{eq:energy-sum}
4\int_{\R^n}\|\nabla\psi\|_2^2
=n+\sum_{j=1}^M\frac{\|v_j\|_2^2}{a_j^2}
\int_\Omega\ddot U\Bigl(\frac{\langle v_j,z\rangle}{a_j}\Bigr)\psi(z)^2\,dz.
\end{equation}
Let us bound each integral on the right. Fix \(j\). For each \(z\), use the orthogonal decomposition:
\[
z=y\,\frac{v_j}{\|v_j\|_2}+w,\quad
y=\frac{\langle z,v_j\rangle}{\|v_j\|_2},\quad w\perp v_j.
\]
Then:
\[
\psi(z)^2=
\frac1Z\,e^{-y^2/2}\,
g\Bigl(\frac{\|v_j\|_2\,y}{a_j}\Bigr)^2
\cdot
e^{-\|w\|_2^2/2}
\prod_{k\ne j}
g\!\left(
\frac{y\left\langle v_k,\frac{v_j}{\|v_j\|_2}\right\rangle+\langle v_k,w\rangle}{a_k}
\right)^2.
\]
Let \(z\sim\psi^2\). Integrating over \(w\in v_j^\perp\), we find the density of \(y\):
\[
p_j(y)
=
\frac1Z\,e^{-y^2/2}\,
g\Bigl(\frac{\|v_j\|_2\,y}{a_j}\Bigr)^2
\int_{v_j^\perp}
e^{-\|w\|_2^2/2}
\prod_{k\ne j}
g\!\left(
\frac{y\left\langle v_k,\frac{v_j}{\|v_j\|_2}\right\rangle+\langle v_k,w\rangle}{a_k}
\right)^2
dw.
\]
Set:
\[
H(y,w):=e^{-\|w\|_2^2/2}
\prod_{k\ne j}
g\!\left(
\frac{y\left\langle v_k,\frac{v_j}{\|v_j\|_2}\right\rangle+\langle v_k,w\rangle}{a_k}
\right)^2,
\qquad
h(y):=\int_{v_j^\perp}H(y,w)\,dw.
\]
Then:
\[
p_j(y)=\frac1Z\,e^{-y^2/2}\,g\Bigl(\frac{\|v_j\|_2\,y}{a_j}\Bigr)^2h(y).
\]
\(g\) is log-concave on \(\R\); hence, each factor of \(H\) is log-concave in the joint variable \((y,w)\).
Thus, \(H(y,w)\) is log-concave. By Pr\'ekopa's theorem (eg, \cite{SaumWell}), \(h(y)\) is also log-concave. 
\smallskip%
 
\(g\) is even; hence, \(H(y,w)=H(-y,-w)\). Thus:
\[
h(y)=\int_{v_j^\perp}H(y,w)\,dw=
\int_{v_j^\perp}H(-y,-w)\,dw=
\int_{v_j^\perp}H(-y,w')\,dw'=h(-y),
\]
and \(h(y)\) is also even.
\smallskip%
 
An even log-concave function is non-increasing on \([0,\infty)\). Therefore:
\begin{equation}\label{h-monotone}
\forall s,t\ge0\qquad h(s+t)\le h(s).
\end{equation}
Let us bound the \(j\)-th integral in \eqref{eq:energy-sum}. We compute:
\[
\int_\Omega\ddot U\Bigl(\frac{\langle v_j,z\rangle}{a_j}\Bigr)\psi^2
=\E_y\,\ddot U\Bigl(\frac{\|v_j\|_2\,y}{a_j}\Bigr)
=\frac1Z\int_{\R}\bigl(g^2\ddot U\bigr)\Bigl(\frac{\|v_j\|_2\,y}{a_j}\Bigr)e^{-y^2/2}h(y)\,dy,
\]
where the integrand is set to zero for \(|y|\ge a_j/\|v_j\|_2\). By \eqref{eq:gU}, \(g^2\ddot U\le C\). Also, \(\ddot U=0\) on \((-1/2,1/2)\). Hence:
\[
\int_\Omega\ddot U\Bigl(\frac{\langle v_j,z\rangle}{a_j}\Bigr)\psi^2
\le C\int_{t<|y|<2t}q(y)\,dy,
\qquad\text{where }\ t:=\frac{a_j}{2\|v_j\|_2},\quad q(y):=\frac1Z\,e^{-y^2/2}h(y).
\]
From \eqref{h-monotone}:
\[
\forall s \ge 0\quad 
q(s+t) = \frac1Z\,e^{-(s+t)^2/2}h(s+t) \le 
\frac1Z\, e^{-s^2/2} e^{-t^2/2} h(s) = 
e^{-t^2/2} q(s).
\]
Hence:
\[
\int_t^{2t}q(y)\,dy\le e^{-t^2/2}\int_0^{t}q(y)\,dy.
\]
\(q\) is even; hence:
\[
\int_{t<|y|<2t}q(y)\,dy = 2 \int_{t<y<2t}q(y)\,dy.
\]
Therefore:
\[
\int_\Omega\ddot U\Bigl(\frac{\langle v_j,z\rangle}{a_j}\Bigr)\psi^2 \le 2 C e^{-t^2/2}\int_0^{t}q(y)\,dy.
\]
For \(|y|\le t\), \(g(\|v_j\|_2\,y/a_j)=1\); hence: \(q(y)=p_j(y)\). Thus:
\[
\int_{0}^{t} q(y)\,dy=
\int_{0}^{t} p_j(y)\,dy = \mathbb P\bigl(0 \le y \le t\bigr)\le1.
\]
Consequently,
\[
\int_\Omega\ddot U\Bigl(\frac{\langle v_j,z\rangle}{a_j}\Bigr)\psi^2
\le 2C e^{-t^2/2}
=2C \exp\Bigl(-\frac{a_j^2}{8\|v_j\|_2^2}\Bigr),
\]
and the proof follows. \qed
\end{proof}

\section{Chaining}\label{chaining}
Let \(X=\{p_1,\ldots,p_n\}\subset\R^2\). 
A \emph{halfplane trace} is a set \(X\cap P\), where \(P\subset\R^2\) is a halfplane. Let \(\mathcal H(X)\) denote the family of halfplane traces. For \(F\subset X\), let \(\one_F\in\R^n\) be its indicator vector: 
\[
(\one_F)_i = \begin{cases}
1, & \text{if \(p_i\in F\)},\\
0, & \text{if \(p_i\notin F\)}.
\end{cases}
\]
Thus:
\[
\langle\one_F,z\rangle=\sum_{p_i\in F}z_i\quad \text{for \(z\in\R^n\)}.
\]
For \(D>0\), set:
\[
K_D(X):=\bigl\{z\in\R^n:\ |\langle\one_F,z\rangle|\le D\ \text{for all }F\in\mathcal H(X)\bigr\}.
\]
\(K_D(X)\) is convex, bounded, and centrally symmetric. It is also an intersection of 
slabs. However, we do not apply Proposition~\ref{slabs} to these slabs directly. Instead, we write each \(\one_F\) as a sum of short increments and impose slabs on the increments.
\smallskip%

We use the following known packing result. See \cite[\S\,5.3]{Mat} 
and \cite{ChazWelzl}.
\begin{theorem}\label{thm:packing}
For \(\delta>0\), a family \(\mathcal F\subset\mathcal H(X)\) is \emph{\(\delta\)-separated} if \(|F\triangle F'|>\delta\) for all distinct \(F,F'\in\mathcal F\). There is an absolute constant \(C_0\) such that for every \(0<\delta\le n\) and every \(\delta\)-separated \(\mathcal F\subset\mathcal H(X)\):
\[
|\mathcal F|\le C_0\,(n/\delta)^2.
\]
\end{theorem}
For \(\delta<1\), Theorem~\ref{thm:packing} is just the bound \(|\mathcal H(X)|=O(n^2)\).
\smallskip%

\emph{Nets.} Let us use Theorem~\ref{thm:packing} to build a hierarchy of nets in \(\mathcal H(X)\). This greedy construction is well 
known. For \(F,G\in\mathcal H(X)\), set:
\[
d(F,G):=|F\triangle G|.
\]
Then:
\[
\|\one_F-\one_G\|_2^2=d(F,G).
\]
Set:
\[
L:=\lceil\log_2n\rceil+1,\qquad
\delta_i:=n2^{-i}\qquad 0\le i\le L.
\]
Set \(\mathcal F_0:=\{\varnothing\}\). For \(1\le i\le L\), let \(\mathcal F_i\subset\mathcal H(X)\) be a maximal \(\delta_i\)-separated family. By Theorem~\ref{thm:packing}, \(|\mathcal F_i|\le C_0\,4^i\).
\smallskip%

By maximality, every \(G\in\mathcal H(X)\) satisfies \(d(G,F)\le\delta_i\) for some \(F\in\mathcal F_i\). Hence, each \(F\in\mathcal F_i\), \(i\ge2\), has a \emph{parent} \(p(F)\in\mathcal F_{i-1}\) with \(d(F,p(F))\le\delta_{i-1}\). For \(F\in\mathcal F_1\), set \(p(F):=\varnothing\); then \(d(F,p(F))=|F|\le\delta_0\). In all cases:
\[
d(F,p(F))\le s_i:=n2^{1-i}\qquad   \forall F\in\mathcal F_i.
\]
Also, \(\delta_L\le1/2\). Since 
\(d\) is integer-valued, \(\mathcal F_L=\mathcal H(X)\).
\smallskip%

\emph{Increments.} For \(1\le i\le L\), set:
\[
\mathcal V_i:=\bigl\{\one_F-\one_{p(F)}:\ F\in\mathcal F_i\bigr\} - \left\{0\right\}.
\]
Fix \(H\in\mathcal H(X)\). Since \(\mathcal F_L=\mathcal H(X)\), we may set \(F_L:=H\). Then set:
\[
F_{i-1}:=p(F_i)\quad 
\text{for \(i = L, \ldots, 1\).}
\]
The increments telescope:
\begin{equation}\label{eq:telescope}
\one_H=\sum_{i=1}^L\bigl(\one_{F_i}-\one_{F_{i-1}}\bigr).
\end{equation}
Along with \eqref{eq:telescope}, below we use the following properties of the increments:
\begin{equation}\label{eq:increments}
|\mathcal V_i|\le C_0\,4^i,\qquad
\|v\|_2^2\le s_i\quad \forall v\in\mathcal V_i.
\end{equation}
\begin{proposition}\label{prop:halfplanes}
For each \(D > 0\), there 
exists \(\psi_D\in C_c^\infty(\R^n)\) such that:
\[
\|\psi_D\|_{L^2(\R^n)}=1,\quad
\operatorname{supp}\psi_D\subset K_{D}(X),
\]
and such that:
\begin{equation}\label{eq:halfplane-energy}
\int_{\R^n}\|\nabla\psi_D\|_2^2\le C_1\,\frac{n^{3/2}}{D^2},
\end{equation}
where \(C_1 > 0\) is an absolute constant.
\end{proposition}
\begin{proof}
Choose \(k_0\) so that:
\[
\sqrt n\le2^{k_0}<2\sqrt n.
\]
One may take \(k_0:=\lceil\frac12\log_2n\rceil\). Set:
\[
b_i:=n^{1/4}\,2^{-|i-k_0|/4},
\]
and set:
\[
K':=\bigl\{z\in\R^n:\ |\langle v,z\rangle|\le b_i\ \text{for all }v\in\mathcal V_i\text{ and all }i\bigr\}.
\]
By \eqref{eq:telescope}, for \(z\in K'\) and \(H\in\mathcal H(X)\):
\[
|\langle\one_H,z\rangle| \le 
\sum_{i=1}^L
|\langle\one_{F_i}-\one_{F_{i-1}},z\rangle| \le\sum_{i=1}^Lb_i,
\]
since each \(\one_{F_i}-\one_{F_{i-1}}\) is either 0 or in \(\mathcal V_i\). Consequently,
\[
|\langle\one_H,z\rangle| \le Bn^{1/4},\qquad
B:=\sum_{i\in\mathbb Z}2^{-|i|/4}.
\]
Hence:
\begin{equation}\label{eq:chain-support}
K'\subset K_{Bn^{1/4}}(X).
\end{equation}
In particular, \(K'\) is bounded.
\smallskip%

\emph{Fisher cost.} Apply Proposition~\ref{slabs} to \(K'\): the vectors are \(v\in\mathcal V_i\), and the width of \(v\in\mathcal V_i\) is \(b_i\). Let \(\psi\) be the resulting function. By \eqref{eq:slab-fisher},
\[
\int_{\R^n}\|\nabla\psi\|_2^2\le
n+O\Bigl(\sum_{i=1}^L\sum_{v\in\mathcal V_i}f_i\bigl(\|v\|_2^2\bigr)\Bigr),
\qquad\text{where }\ f_i(x):=\frac{x}{b_i^2}\exp\Bigl(-\frac{b_i^2}{8x}\Bigr).
\]
Each \(f_i\) is increasing on \((0,\infty)\). By \eqref{eq:increments}, \(\|v\|_2^2\le s_i\) for \(v\in\mathcal V_i\); hence, \(f_i(\|v\|_2^2)\le f_i(s_i)\). Also by \eqref{eq:increments}, the inner sum has at most \(C_0\,4^i\) terms. Therefore:
\[
\int_{\R^n}\|\nabla\psi\|_2^2\le
n+O\Bigl(\sum_{i=1}^L4^i\,\frac{s_i}{b_i^2}\exp\Bigl(-\frac{b_i^2}{8s_i}\Bigr)\Bigr).
\]
Let us bound 
the \(i\)-th term. Write: 
\[
j:=i-k_0.
\] 
Since \(2^{k_0}<2\sqrt n\), 
\[
4^{i} = 4^{j} 4^{k_0} = O(n\,4^j).
\]
Since \(2^{k_0}\ge\sqrt n\), 
\[
s_i= n 2^{1-i} = 2n2^{-k_0}\,2^{-j}\le2\sqrt n\,2^{-j}.
\]
Since \(b_i^2=\sqrt n\,2^{-|j|/2}\),
\[
4^i\,\frac{s_i}{b_i^2}=O\bigl(n\,2^{j+|j|/2}\bigr),\qquad
\frac{b_i^2}{8s_i}\ge\frac{2^{j-|j|/2}}{16}.
\]
\emph{Case \(j \le 0\).} In this case, 
\[
4^i\,\frac{s_i}{b_i^2} = 
O\bigl(n\,2^{j+|j|/2}\bigr) = 
O\bigl(n\,2^{-|j|/2}\bigr).
\]
We bound the \(i\)-th term as follows:
\[
4^i\,\frac{s_i}{b_i^2}\exp\Bigl(-\frac{b_i^2}{8s_i}\Bigr) \le 
4^i\,\frac{s_i}{b_i^2} = O\bigl(n\,2^{-|j|/2}\bigr).
\]
\emph{Case \(j > 0\).} In this case, 
\[
4^i\,\frac{s_i}{b_i^2} = 
O\bigl(n\,2^{j+|j|/2}\bigr) = 
O\bigl(n\,2^{3j/2}\bigr),\qquad 
\frac{b_i^2}{8s_i}\ge\frac{2^{j-|j|/2}}{16} = \frac{2^{j/2}}{16}.
\]
We bound the \(i\)-th term as follows:
\[
4^i\,\frac{s_i}{b_i^2}\exp\Bigl(-\frac{b_i^2}{8s_i}\Bigr) \le 
O\bigl(n\,2^{3j/2}\bigr) \cdot 
\exp\bigl(-\frac{2^{j/2}}{16}\bigr) = O\bigl(n\,2^{3j/2}\bigr) \cdot O(2^{-2j}) = O\bigl(n\,2^{-j/2}\bigr).
\]
In both cases, the \(i\)-th term is \(O(n\,2^{-|j|/2})\). Therefore:
\[
\sum_{i=1}^L4^i\,\frac{s_i}{b_i^2}\exp\Bigl(-\frac{b_i^2}{8s_i}\Bigr) \le 
O\bigl(
\sum_{j\in\mathbb Z} n\,2^{-|j|/2}
\bigr) = O(n).
\]
Consequently,
\[
\int_{\R^n}\|\nabla\psi\|_2^2\le
n+O\Bigl(\sum_{i=1}^L4^i\,\frac{s_i}{b_i^2}\exp\Bigl(-\frac{b_i^2}{8s_i}\Bigr)\Bigr) = n + O(n) = O(n).
\]
\emph{Scaling.} Now, we rescale \(\psi\). Set:
\[
r:=D/(Bn^{1/4}),\quad 
\psi_D(z):=r^{-n/2}\psi(z/r).
\]
Then \(\|\psi_D\|_{L^2(\R^n)}=1\). Also, \(\operatorname{supp}\psi_D=r\operatorname{supp}\psi\). By \eqref{eq:chain-support},
\[
\operatorname{supp}\psi_D\subset rK_{Bn^{1/4}}(X)=K_{D}(X).
\]
Finally, 
\[
\int\|\nabla\psi_D\|_2^2=r^{-2}\int\|\nabla\psi\|_2^2,\quad 
\int_{\R^n}\|\nabla\psi_D\|_2^2=O\bigl(r^{-2}n\bigr)=O\Bigl(\frac{n^{3/2}}{D^2}\Bigr),
\]
and the proof follows. \qed
\end{proof}

\section{Prefix halfplanes}\label{sec:prefix}
Recall the setting of \S\,\ref{chaining}. 
Let again \(X=\{p_1,\ldots,p_n\}\subset\R^2\) be a set of points, \(\mathcal H(X)\) be the family of halfplane traces. For a subset \(F\subset X\), let \(\one_F\in\R^n\) be its 
indicator vector. Set again:
\[
K_D(X)=\bigl\{z\in\R^n:\ |\langle\one_F,z\rangle|\le D\ \text{for all }F\in\mathcal H(X)\bigr\}.
\]
Now, let \(X=\{p_1,\ldots,p_n\}\subset\R^2\) be an \emph{ordered} set of points. For \(0\le k\le n\), set:
\[
X_k:=\{p_1,\ldots,p_k\}.
\]
For \(\chi\in\R^n\) and \(F\subset X\), write:
\[
\chi(F):=\langle\one_F,\chi\rangle=\sum_{p_i\in F}\chi_i.
\]
\begin{theorem}\label{thm:prefix}
Every \(X=\{p_1,\ldots,p_n\}\) has a coloring \(\chi\in\{-1,1\}^n\) with:
\[
|\chi(F\cap X_k)|=O(n^{1/4})\qquad\forall F\in\mathcal H(X),\ \ 0\le k\le n.
\]
\end{theorem}
The proof relies on the following result of \cite{Versh-Sm-3}. 
\begin{theorem}\label{thm:SV}
Let \(K\subset\R^n\) be convex and centrally symmetric. Let \(\psi\in C_c^\infty(\R^n)\) be such that:
\[
\|\psi\|_{L^2(\R^n)}=1,\qquad 
\operatorname{supp}\psi\subset K.
\]
Let \(v_1,\ldots,v_m\in\R^n\) be any vectors, and let \(\e_1,\ldots,\e_m\) be iid fair signs. Then there is an online algorithm that discards some terms \(\e_iv_i\) so that, for each \(k\):
\[
\sum_{i\le k\ \mathrm{accepted}}\e_iv_i\in2K
\quad\text{deterministically},
\]
and:
\[
\E\,\#\{\text{\rm discarded}\}\le\sum_{i=1}^m\bigl\|\langle\nabla\psi,v_i\rangle\bigr\|_{L^2(\R^n)}.
\]
\end{theorem}
First, we color most of the points.
\begin{lemma}\label{lem:partial}
Every \(X=\{p_1,\ldots,p_n\}\) has a partial coloring \(\chi\in\{-1,0,1\}^n\) with at most \(n/4\) zeros and with:
\[
|\chi(F\cap X_k)|=O(n^{1/4})\qquad\forall F\in\mathcal H(X),\ \ 0\le k\le n.
\]
\end{lemma}
\begin{proof}
Fix \(D>0\), and choose 
it later. Apply Theorem~\ref{thm:SV} to \(K=K_D(X)\): choose \(\psi=\psi_D\) from Proposition~\ref{prop:halfplanes}, and \(v_i:=e_i\), \(i=1,\ldots,n\). Set:
\[
\chi_i:=\begin{cases}
\e_i, & \text{if \(\e_ie_i\) is accepted},\\
0, & \text{if \(\e_ie_i\) is discarded}.
\end{cases}
\]
Then:
\[
\forall k\quad \sum_{i\le k}\chi_ie_i\in2K_D(X).
\]
Consequently,
\begin{equation}\label{eq:partial-prefix}
|\chi(F\cap X_k)|\le2D\qquad\forall F\in\mathcal H(X),\ \ 0\le k\le n.
\end{equation}
By Theorem~\ref{thm:SV},
\[
\E\,\#\{\chi_i=0\}
\le\sum_{i=1}^n\|\partial_i\psi_D\|_{L^2(\R^n)}.
\]
By Cauchy-Schwarz,
\[
\E\,\#\{\chi_i=0\}
\le \sqrt n\,\Bigl(\int_{\R^n}\|\nabla\psi_D\|_2^2\Bigr)^{1/2}.
\]
From \eqref{eq:halfplane-energy}:
\[
\E\,\#\{\chi_i=0\} =O\Bigl(\frac{n^{5/4}}{D}\Bigr).
\]
Choose \(D=O(n^{1/4})\) so that 
\(\E\,\#\{\chi_i=0\}\le n/4\). Then there exists some \(\chi\) with at most \(n/4\) zeros, and the proof follows. \qed
\end{proof}
\smallskip%

\begin{proof}[Proof of Theorem~\ref{thm:prefix}]
The following iteration is standard; see \cite{Spencer-6}. 
\smallskip%

We iterate Lemma~\ref{lem:partial}, as in 
\cite[\S\,2]{Versh-Sm-4}. First, apply it to \(X\). 
Then apply it to the points left uncolored, in the induced order, and so on until all points are colored. Round \(r\) is applied to at most \(n4^{-r}\) points. A halfplane trace of \(X\), restricted to these points, is a halfplane trace of them. A prefix of \(X\), restricted to these points, 
is a prefix of them. Let \(\chi\) be the 
final coloring. Then \(\chi(F\cap X_k)\) splits into one term per round. By Lemma~\ref{lem:partial}, the term from round \(r\) is \(O\bigl((n4^{-r})^{1/4}\bigr)\); therefore:
\[
|\chi(F\cap X_k)|\le\sum_{r\ge0}O\bigl((n4^{-r})^{1/4}\bigr)=O(n^{1/4}),
\]
and the proof follows. \qed
\end{proof}

\section{Triangles}\label{sec:triangles}
Let \(X=\{p_1,\ldots,p_n\}\subset\R^2\) be a set of points. For \(\chi\in\R^n\) and \(F\subset X\), write again:
\[
\chi(F)=\sum_{p_i\in F}\chi_i.
\]
\begin{theorem}
Every \(X=\{p_1,\ldots,p_n\}\) has a coloring \(\chi\in\{-1,1\}^n\) with:
\[
|\chi(X\cap T)|=O(n^{1/4})\qquad\text{for all triangles }T.
\]
\end{theorem}
\begin{proof}
The proof reduces triangles to corners.  For a nonvertical line \(\ell\colon y=\ell(x)\) and \(c\in\R\), set:
\[
\Gamma(\ell,c):=\bigl\{(x,y)\in\R^2:\ x\le c,\ \ y\le\ell(x)\bigr\}.
\]
Order \(X\) by the first coordinate. Let \(\chi\) be the coloring from Theorem~\ref{thm:prefix} for this order.
\colorlet{rgn}{blue!22}
\colorlet{rgnedge}{blue!55!black}
\tikzset{
  pic/.style={scale=0.55,baseline={([yshift=-.5ex]current bounding box.center)}},
  fade/.style={top color=rgn,bottom color=rgn!0,draw=none},
  solid region/.style={fill=rgn,draw=rgnedge,line width=0.6pt,line join=round},
  edge/.style={draw=rgnedge,line width=0.6pt,line join=round},
  ghost/.style={black!40,thin},
  guide/.style={black!45,densely dashed,thin},
  lab/.style={font=\scriptsize,inner sep=1.5pt}}
\newcommand{\picbox}{\useasboundingbox (-0.1,-0.55) rectangle (4.45,3);}
\newcommand{\guide}[2]{\draw[guide] (#1,0)--(#1,3);\node[lab,anchor=north] at (#1,0) {\(#2\)};}
\newcommand{\ghostline}{\draw[ghost] (0,0.9)--(4,2.5);\node[lab,anchor=west] at (4,2.5) {\(\ell\)};}
\newcommand{\op}[1]{\;#1\;}
\medskip%

\emph{Corners.} \(X\cap\Gamma(\ell,c)\) is a halfplane trace intersected with a prefix:
\[
\begin{tikzpicture}[pic]\picbox
\shade[fade] (0,0)--(2.6,0)--(2.6,1.94)--(0,0.9)--cycle;
\ghostline\guide{2.6}{c}
\draw[edge] (0,0.9)--(2.6,1.94)--(2.6,0);
\end{tikzpicture}
\op{=}
\begin{tikzpicture}[pic]\picbox
\shade[fade] (0,0)--(4,0)--(4,2.5)--(0,0.9)--cycle;
\ghostline
\draw[edge] (0,0.9)--(4,2.5);
\end{tikzpicture}
\op{\cap}
\begin{tikzpicture}[pic]\picbox
\shade[fade] (0,0)--(2.6,0)--(2.6,3)--(0,3)--cycle;
\guide{2.6}{c}
\draw[edge] (2.6,3)--(2.6,0);
\end{tikzpicture}
\]
By Theorem~\ref{thm:prefix},
\begin{equation}\label{eq:corner}
|\chi(X\cap \Gamma(\ell,c))|=O(n^{1/4}).
\end{equation}
\smallskip%

\emph{Interfixes.} Consider the domain below 
\(\ell\) and between the vertical lines \(x=a\) and \(x=b\); it is a difference of two corners:
\[
\begin{tikzpicture}[pic]\picbox
\shade[fade] (1,0)--(3,0)--(3,2.1)--(1,1.3)--cycle;
\ghostline\guide{1}{a}\guide{3}{b}
\draw[edge] (1,0)--(1,1.3)--(3,2.1)--(3,0);
\end{tikzpicture}
\op{=}
\begin{tikzpicture}[pic]\picbox
\shade[fade] (0,0)--(3,0)--(3,2.1)--(0,0.9)--cycle;
\ghostline\guide{3}{b}
\draw[edge] (0,0.9)--(3,2.1)--(3,0);
\end{tikzpicture}
\op{-}
\begin{tikzpicture}[pic]\picbox
\shade[fade] (0,0)--(1,0)--(1,1.3)--(0,0.9)--cycle;
\ghostline\guide{1}{a}
\draw[edge] (0,0.9)--(1,1.3)--(1,0);
\end{tikzpicture}
\]
Its trace is a halfplane trace intersected with an interfix \(X_k-X_j\). By \eqref{eq:corner}, its \(\chi\) is \(O(n^{1/4})\).
\smallskip%

\emph{Triangles with a vertical side.} Such a triangle is a difference of two domains of the previous kind:
\[
\begin{tikzpicture}[pic]\picbox
\draw[ghost] (0.3,2.985)--(4,0.95); \draw[ghost] (0.3,0.42)--(4,1.9);
\node[lab,anchor=east] at (0.3,2.985) {\(\ell_2\)}; \node[lab,anchor=east] at (0.3,0.42) {\(\ell_1\)};
\guide{1}{a}\guide{3}{b}
\draw[solid region] (1,0.7)--(1,2.6)--(3,1.5)--cycle;
\end{tikzpicture}
\op{=}
\begin{tikzpicture}[pic]\picbox
\shade[fade] (1,0)--(3,0)--(3,1.5)--(1,2.6)--cycle;
\draw[ghost] (0.3,2.985)--(4,0.95);
\guide{1}{a}\guide{3}{b}
\draw[edge] (1,0)--(1,2.6)--(3,1.5)--(3,0);
\node[lab,anchor=east] at (0.3,2.985) {\(\ell_2\)};
\end{tikzpicture}
\op{-}
\begin{tikzpicture}[pic]\picbox
\shade[fade] (1,0)--(3,0)--(3,1.5)--(1,0.7)--cycle;
\draw[ghost] (0.3,0.42)--(4,1.9);
\guide{1}{a}\guide{3}{b}
\draw[edge] (1,0)--(1,0.7)--(3,1.5)--(3,0);
\node[lab,anchor=east] at (0.3,0.42) {\(\ell_1\)};
\end{tikzpicture}
\]
For such a triangle, its \(\chi\) is again \(O(n^{1/4})\). 
\smallskip%

\emph{Triangles.} Every triangle \(T\) is a union of 
two triangles with vertical sides. Perturbing \(T\) slightly, we may assume that no point of \(X\) lies on the boundaries in the pictures. Hence, 
\(\chi(X\cap T)\) is the sum of two terms, each \(O(n^{1/4})\). \qed
\end{proof}


\bibliographystyle{plain}
\bibliography{ref}

\end{document}


